\chapter{MATERYAL VE YÖNTEM}

Şekil ve çizelgeler metinde ilk kez anıldıktan sonra uygun en yakın yere yerleştirilmelidir. Örnek olarak Şekil~\ref{fig:ornek-sekil} ve Çizelge~\ref{tab:ornek-cizelge} verilmiştir.

\begin{figure}[htbp]
\centering
\fbox{\rule{0pt}{3cm}\rule{6cm}{0pt}}
\caption{Tek satır olması durumunda ortalı olacak.}
\label{fig:ornek-sekil}
\end{figure}

\begin{table}[htbp]
\caption{Tek satırlı ve kolonlar ortalanmış çizelge.}
\label{tab:ornek-cizelge}
\centering
\begin{tabular}{c c c c}
\toprule\toprule
Kolon A & Kolon B & Kolon C & Kolon D\\
\midrule
Satır A & Satır A & Satır A & Satır A\\
Satır B & Satır B & Satır B & Satır B\\
Satır C & Satır C & Satır C & Satır C\\
\bottomrule
\end{tabular}
\end{table}

\section{Denklemler}
Denklemler 12 punto yazılır ve denklem numaraları sağa dayalı olur. Örneğin Denklem~\ref{eq:ornek} aşağıda verilmiştir.

\begin{equation}
E = mc^2
\label{eq:ornek}
\end{equation}

Parametreler denklemin altında açıklanmalıdır.
